\documentclass[11pt]{article}
\ifdefined\pdfinfoomitdate\pdfinfoomitdate=1\fi
\ifdefined\pdftrailerid\pdftrailerid{}\fi
\usepackage[T1]{fontenc}
\usepackage{lmodern}
\usepackage{amsmath,amssymb,amsthm,booktabs}
\usepackage[a4paper,margin=25mm]{geometry}
\usepackage[colorlinks=true,linkcolor=blue,urlcolor=blue]{hyperref}
\setlength{\emergencystretch}{2em}
\allowdisplaybreaks
\title{Integrality of \(\Gamma_H(G)\) forces \(H\) to be normal}
\author{}
\date{}
\begin{document}
\maketitle

\section*{Statement and source problem}

Let \(G\) be a finite group of order \(n\geq 5\), and let \(H\) be a
proper, nontrivial subgroup of \(G\). Following Definition 1.4 and the Section 2 subgroup convention of Poddar,
Biswas, and Das, set
\[
 S_H=\{(g,1),(1,g),(g,g):g\in G\setminus H\}
\]
and
\[
 \Gamma_H(G)=\operatorname{Cay}(G\times G,S_H).
\]
Thus the vertices are the elements of \(G\times G\), and distinct vertices
\(u,v\) are adjacent precisely when \(uv^{-1}\in S_H\).

The source proves in Theorem 3.1 that \(H\lhd G\) implies that
\(\Gamma_H(G)\) is integral. Its first bullet in Section 6, ``Conclusion and
Open Issues,'' asks whether normality is also necessary. That bullet is
unnumbered in the source; it is not assigned a conjecture number.

Integral means that every adjacency eigenvalue is an integer. We prove the converse.

\textbf{Theorem.} The graph \(\Gamma_H(G)\) is integral if and only if
\(H\lhd G\).

\section*{Proof idea}

Assume that \(H\) is not normal. The permutation representation of \(G\) on
the cosets \(G/H\) must then contain an irreducible representation
\((\rho,V)\) whose \(H\)-fixed subspace is neither zero nor all of \(V\).
This representation detects the failure of normality.

The adjacency operator of a Cayley graph splits into blocks indexed by the
irreducible representations of the underlying group. We use the block
\(\rho\otimes\rho^*\) of \(G\times G\). After identifying
\(V\otimes V^*\) with \(\operatorname{End}(V)\), that block preserves the
two-dimensional space spanned by the orthogonal projection onto \(V^H\)
and its complementary projection. The characteristic polynomial on this
plane has a discriminant equal to a rational square factor times a positive
nonsquare integer. It therefore contributes two irrational adjacency
eigenvalues.

\section*{1. An irreducible representation detecting nonnormality}

We prove the contrapositive, so suppose that \(H\) is not normal. Use the
source notation
\[
 n=|G|,\qquad k=|H|,\qquad \ell=[G:H],
\]
so that \(n=k\ell\).

Consider the complex permutation module on the left cosets,
\[
 \mathbb C[G/H]=\operatorname{Ind}_H^G(\mathbf 1_H).
\]
For \((\sigma,W)\in\operatorname{IRR}(G)\), Frobenius reciprocity says that
the multiplicity of \(\sigma\) in this module is
\[
 \dim W^H.
\]
Suppose that every irreducible constituent of \(\mathbb C[G/H]\) were
fixed pointwise by \(H\). Then \(H\) would act trivially on the entire
permutation module and hence on every coset. Thus, for all \(a\in H\) and
\(g\in G\),
\[
 agH=gH.
\]
Equivalently, \(g^{-1}ag\in H\). It would follow that
\(g^{-1}Hg\subseteq H\); equality of orders would give
\(g^{-1}Hg=H\) for every \(g\), contrary to the nonnormality of \(H\).

Consequently there is an irreducible representation
\((\rho,V)\in\operatorname{IRR}(G)\) for which
\[
 0<r:=\dim V^H<d:=\deg(\rho)=\dim V. \tag{1}
\]
In particular, \(\rho\) is nontrivial and \(d\geq2\).

Taking dimensions in the irreducible decomposition of
\(\mathbb C[G/H]\) gives
\[
 \ell=\sum_{(\sigma,W)\in\operatorname{IRR}(G)}
       (\dim W)(\dim W^H).
\]
The trivial representation contributes \(1\), while the chosen
representation contributes \(dr\). Every other term is nonnegative, so
\[
 dr\leq \ell-1,\qquad\text{hence}\qquad \ell\geq dr+1. \tag{2}
\]

Choose a \(G\)-invariant Hermitian inner product on \(V\), making \(\rho\)
unitary. In the source notation for sums of representation matrices, put
\[
 P=P_H:=\frac{1}{k}\rho(H)
       =\frac{1}{k}\sum_{a\in H}\rho(a),
 \qquad Q=I-P.
\]
Then \(P\) is the orthogonal projection onto \(V^H\), and therefore
\[
 P^2=P=P^*,\qquad \operatorname{Tr}(P)=r. \tag{3}
\]
By (1), both \(P\) and \(Q\) are nonzero.

\section*{2. A two-dimensional adjacency block}

The contragredient (dual) representation \(\rho^*\) is irreducible, so
\[
 \Omega:=\rho\otimes\rho^*\in\operatorname{IRR}(G\times G).
\]
Identify \(V\otimes V^*\) with \(\operatorname{End}(V)\). Under this
identification,
\[
 \Omega(a,b)X=\rho(a)X\rho(b)^{-1}
 \qquad(a,b\in G,\ X\in\operatorname{End}(V)). \tag{4}
\]

The regular representation of \(G\times G\) contains every irreducible
representation. Hence the adjacency matrix of \(\Gamma_H(G)\) contains the
Fourier block
\[
 B:=\Omega(S_H)=\sum_{s\in S_H}\Omega(s).
\]
Explicitly,
\[
 B(X)=\sum_{g\in G\setminus H}
 \bigl(\rho(g)X+X\rho(g)^{-1}+\rho(g)X\rho(g)^{-1}\bigr). \tag{5}
\]
Every eigenvalue of \(B\) is therefore an adjacency eigenvalue of
\(\Gamma_H(G)\).

We now evaluate \(B\) on \(P\) and \(Q\). Since \(\rho\) is nontrivial and
irreducible,
\[
 \sum_{g\in G}\rho(g)=0. \tag{6}
\]
Schur averaging on \(\operatorname{End}(V)\) gives
\[
 \sum_{g\in G}\rho(g)X\rho(g)^{-1}
   =\frac{n}{d}\operatorname{Tr}(X)I. \tag{7}
\]
Also, inversion permutes both \(H\) and \(G\setminus H\), so
\[
 \sum_{g\in G\setminus H}\rho(g)
 =\sum_{g\in G\setminus H}\rho(g)^{-1}
 =-kP. \tag{8}
\]
Finally, \(P\) and \(Q\) commute with every \(\rho(a)\), \(a\in H\).
Subtracting the \(H\)-conjugation average from (7) and using (8) in (5)
yields
\[
 B(P)=\left(\frac{nr}{d}-3k\right)P+\frac{nr}{d}Q, \tag{9}
\]
and
\[
 B(Q)=\frac{n(d-r)}{d}P+
       \left(\frac{n(d-r)}{d}-k\right)Q. \tag{10}
\]

Thus \(\langle P,Q\rangle\) is a two-dimensional \(B\)-invariant subspace.
Set
\[
 \alpha=\frac{nr}{d},\qquad
 \beta=\frac{n(d-r)}{d}.
\]
In the ordered basis \((P,Q)\), the restricted operator has matrix
\[
 \begin{pmatrix}
   \alpha-3k & \beta\\
   \alpha    & \beta-k
 \end{pmatrix}. \tag{11}
\]
Its trace and determinant are
\[
 \operatorname{Tr}=n-4k,\qquad
 \det=3k^2-3nk+\frac{2nkr}{d}. \tag{12}
\]
Its discriminant is therefore
\[
 \begin{aligned}
 \Delta
   &=(n-4k)^2
     -4\left(3k^2-3nk+\frac{2nkr}{d}\right)\\
   &=(n+2k)^2-\frac{8nkr}{d}\\
   &=\left(\frac{k}{d}\right)^2 M,
 \end{aligned} \tag{13}
\]
where, since \(n=k\ell\),
\[
 M=d^2(\ell+2)^2-8\ell dr. \tag{14}
\]

\section*{3. The integer \(M\) is not a square}

Define
\[
 N=d(\ell+2)-4r.
\]
A direct expansion of (14) gives
\[
 M=N^2+16r(d-r). \tag{15}
\]
By (1), \(M>N^2\). Moreover, (2) implies
\[
 N\geq d(dr+3)-4r
   =r(d^2-4)+3d>0. \tag{16}
\]

Except at two smallest parameter triples, \(M\) lies strictly between
\(N^2\) and \((N+1)^2\); the two exceptions will be bracketed separately.
For the general cases, the relevant gap is
\[
 (N+1)^2-M
   =2d(\ell+2)-8r+1-16r(d-r). \tag{17}
\]
For fixed \(d,r\), this expression increases by \(2d\) when \(\ell\)
increases by \(1\). At the smallest value allowed by (2),
\(\ell=dr+1\), it becomes
\[
 F(d,r)=2rd(d-8)+16r^2-8r+6d+1. \tag{18}
\]

For \(d\geq8\), all three terms
\[
 2rd(d-8),\qquad 16r^2-8r,\qquad 6d+1
\]
are nonnegative, and the last two are positive. For \(4\leq d\leq7\), the
same positivity follows from the following forms:

\begin{center}
\begin{tabular}{ll}
\toprule
\(d\) & \(F(d,r)\) \\
\midrule
\(4\) & \((4r-5)^2\) \\
\(5\) & \(16(r-1)^2-6(r-1)+9\) \\
\(6\) & \(16(r-1)^2+21\) \\
\(7\) & \(16(r-1)^2+10(r-1)+37\) \\
\bottomrule
\end{tabular}
\end{center}

For \(d=5\), put \(t=r-1\geq0\). At \(t=0\) the expression is 9; at integer \(t\geq1\), \(16t^2-6t+9\geq10t+9>0\). These expressions are positive for every integer \(1\leq r\leq d-1\).
The remaining case \(d=3,r=2\) gives \(F(3,2)=7>0\). In all these cases,
(15), (17), and (18) imply
\[
 N^2<M<(N+1)^2. \tag{19}
\]

There are only two exceptional smallest triples.

\begin{itemize}
\item If \((d,r,\ell)=(3,1,4)\), then \(N=14\) and \(M=228\), so
  \[
  15^2=225<228<256=16^2.
  \]
  For every \(\ell\geq5\), the expression in (17) is \(-3+6(\ell-4)>0\).
\item If \((d,r,\ell)=(2,1,3)\), then \(N=6\) and \(M=52\), so
  \[
  7^2=49<52<64=8^2.
  \]
  For every \(\ell\geq4\), the expression in (17) is \(-3+4(\ell-3)>0\).

\end{itemize}
This exhausts \(d\geq2\) and \(1\leq r\leq d-1\). Hence \(M\) is a positive
integer that is not an integer square. An integer that is a square in
\(\mathbb Q\) is already an integer square, so \(\sqrt M\notin\mathbb Q\).

\section*{4. Irrational adjacency eigenvalues}

By (12)--(14), the two eigenvalues of the restriction (11) are
\[
 \lambda_\pm
  =\frac{n-4k\pm (k/d)\sqrt M}{2}. \tag{20}
\]
Both are irrational. Since this restriction belongs to the genuine
\(\Omega=\rho\otimes\rho^*\) adjacency block, both eigenvalues occur in
\(sp(\Gamma_H(G))\). Thus \(\Gamma_H(G)\) is not integral whenever
\(H\) is not normal.

This proves the contrapositive: if \(\Gamma_H(G)\) is integral, then
\(H\lhd G\). Together with Theorem 3.1 of the source, the equivalence follows.
\(\square\)

\section*{Dependencies and scope}

The proof uses complete reducibility for finite groups over \(\mathbb C\),
Frobenius reciprocity, unitary realization of finite-group
representations, the irreducibility of tensor products for direct products,
the regular-representation decomposition of a Cayley adjacency operator,
and Schur averaging. It uses no classification of finite groups and no
finite enumeration. The assumptions on \(G\), \(H\), and \(\Gamma_H(G)\)
are those in Definition 1.4, the subgroup convention of Section 2, and the first open issue of Section 6.

\section*{Source}

S. Poddar, S. Biswas, and A. Das, ``An integral family of quasi-strongly
regular Cayley graphs,'' \href{https://arxiv.org/abs/2511.14496v1}{arXiv:2511.14496v1}.
The graph \(\Gamma_H(G)\) is Definition 1.4, the implication
\(H\lhd G\Rightarrow\Gamma_H(G)\) integral is Theorem 3.1, and the converse
proved above is the first unnumbered open issue in Section 6.
\end{document}
